\documentclass[12pt,a4paper,twoside]{report}
\usepackage[utf8]{inputenc}
\usepackage[T1]{fontenc}
\usepackage[polish]{babel}
\usepackage{times}
\usepackage{indentfirst}
\usepackage{setspace}
\usepackage[margin=2.5cm]{geometry}
\usepackage{fancyhdr}
\usepackage{titlesec}
\usepackage{hyperref}
\usepackage{enumitem}
\usepackage{graphicx}
\usepackage{caption}
\usepackage{float}
\usepackage{url}

% Ustawienia odstępów
\onehalfspacing
\setlength{\parindent}{0.75cm}
\setlength{\parskip}{0pt}

% Formatowanie nagłówków rozdziałów
\titleformat{\chapter}[display]
  {\normalfont\fontsize{16}{18}\bfseries}{\chaptertitlename\ \thechapter}{12pt}{}
\titlespacing*{\chapter}{0pt}{-20pt}{12pt}

\titleformat{\section}
  {\normalfont\fontsize{14}{16}\bfseries}{\thesection}{6pt}{}
\titlespacing*{\section}{0pt}{12pt}{6pt}

\titleformat{\subsection}
  {\normalfont\fontsize{12}{14}\bfseries}{\thesubsection}{6pt}{}
\titlespacing*{\subsection}{0pt}{6pt}{6pt}

% Głębokość numerowania
\setcounter{secnumdepth}{3}
\setcounter{tocdepth}{2}

% Definicje dla spisów
\usepackage{tocloft}
\renewcommand{\cftchapfont}{\normalfont}
\renewcommand{\cftchappagefont}{\normalfont}

% Styl strony - początkowo numeracja wyśrodkowana na dole
\fancypagestyle{plain}{
    \fancyhf{}
    \fancyfoot[C]{\thepage}
    \renewcommand{\headrulewidth}{0pt}
}
% Styl dla stron z treścią rozdziałów - numeracja w prawym dolnym rogu
\fancypagestyle{main}{
    \fancyhf{}
    \fancyfoot[R]{\thepage}
    \renewcommand{\headrulewidth}{0pt}
}

% Przełączanie stylów
\makeatletter
\let\ps@plain\ps@main
\makeatother

\begin{document}

% Strona tytułowa (bez numeracji)
\begin{titlepage}
    \centering
    \vspace*{2cm}
    {\fontsize{18}{22}\selectfont \textbf{AKADEMIA NAUK STOSOWANYCH W NOWYM SĄCZU}\par}
    \vspace{1cm}
    {\fontsize{18}{22}\selectfont \textbf{WYDZIAŁ NAUK INŻYNIERYJNYCH}\par}
    \vspace{2cm}
    {\fontsize{16}{20}\selectfont \textbf{PRACA DYPLOMOWA}\par}
    \vspace{1.5cm}
    {\fontsize{16}{20}\selectfont \textbf{Zdecentralizowany system obiegu dokumentów z wykorzystaniem blockchain i IPFS}\par}
    \vspace{2cm}
    \begin{flushright}
        {\fontsize{14}{18}\selectfont Autor: Mateusz Smaga\par}
        \vspace{0.3cm}
        {\fontsize{12}{15}\selectfont Kierunek: Informatyka Stosowana\par}
        \vspace{0.3cm}
        {\fontsize{12}{15}\selectfont Nr albumu: 33441\par}
    \end{flushright}
    \vspace{2cm}
    {\fontsize{12}{15}\selectfont Promotor: dr Grzegorz Litawa\par}
    \vspace{1cm}
    {\fontsize{12}{15}\selectfont NOWY SĄCZ 2026\par}
\end{titlepage}

% Spis treści (numeracja wyśrodkowana)
\pagestyle{plain}
\tableofcontents
\clearpage

% Przełączamy styl na main (prawy dolny róg) od pierwszej strony rozdziału
\pagestyle{main}

% Rozdział 1

% ==================== ROZDZIAŁ 1: PODSTAWY TEORETYCZNE (ROZSZERZONY) ====================
\chapter{Podstawy teoretyczne systemu}

\section{Istota i architektura technologii Blockchain}

Technologia blockchain, klasyfikowana jako odmiana technologii rozproszonego rejestru
(ang.~\textit{Distributed Ledger Technology}), stanowi fundamentalną zmianę w sposobie
przechowywania i wymiany danych cyfrowych. W odróżnieniu od tradycyjnych, scentralizowanych
baz danych zarządzanych przez pojedynczy podmiot, blockchain funkcjonuje jako sieć
równorzędnych uczestników (węzłów), z których każdy przechowuje pełną lub częściową kopię
historii wszystkich transakcji.

Struktura danych opiera się na sekwencyjnym łączeniu bloków. Każdy blok, oprócz zestawu
transakcji, zawiera w swoim nagłówku znacznik czasu (ang.~\textit{timestamp}) oraz
kryptograficzny skrót poprzedniego bloku. Wykorzystanie jednokierunkowej funkcji skrótu
sprawia, że jakakolwiek próba modyfikacji danych w bloku historycznym spowoduje zmianę jego
skrótu, a w konsekwencji zerwanie łańcucha referencji we wszystkich kolejnych blokach. Taka
konstrukcja zapewnia matematycznie gwarantowaną niezmienność (ang.~\textit{immutability})
zapisanych danych -- raz zatwierdzona informacja staje się trwałym elementem rejestru,
którego nie można usunąć ani sfałszować bez wykrycia przez pozostałych uczestników sieci.

Technologia ta została po raz pierwszy zaimplementowana w praktyce przez anonimowego twórcę
posługującego się pseudonimem Satoshi Nakamoto w kontekście kryptowaluty
Bitcoin~\cite{nakamoto2008bitcoin}. Od tego czasu koncepcja ta ewoluowała, wykraczając
daleko poza zastosowania czysto finansowe. W kontekście niniejszej pracy blockchain
traktowany jest nie jako system płatniczy, lecz jako zdecentralizowana infrastruktura
zaufania -- publiczny rejestr zdolny do przechowywania dowodów istnienia i integralności
dokumentów cyfrowych.

Podstawowe cechy charakterystyczne tej technologii obejmują:
\begin{itemize}
    \item \textbf{Decentralizację i odporność na awarie:} Brak centralnego serwera eliminuje
          pojedynczy punkt awarii (ang.~\textit{Single Point of Failure}). Nawet w przypadku
          wyłączenia znacznej części węzłów sieci dane pozostają dostępne i bezpieczne
          u pozostałych uczestników.
    \item \textbf{Niezmienność i integralność danych:} Dzięki kryptograficznemu powiązaniu
          bloków modyfikacja danych historycznych wymagałaby ponownego wykonania pracy
          obliczeniowej dla wszystkich kolejnych bloków, co jest praktycznie niewykonalne.
    \item \textbf{Transparentność i audytowalność:} W publicznych sieciach blockchain każdy
          może niezależnie zweryfikować poprawność całego łańcucha oraz autentyczność
          pojedynczego wpisu, korzystając z publicznie dostępnych eksploratorów bloków.
    \item \textbf{Odporność na cenzurę:} Transakcja spełniająca reguły protokołu zostanie
          włączona do bloku niezależnie od woli jakiegokolwiek zewnętrznego podmiotu
          czy instytucji.
\end{itemize}

% -----------------------------------------------------------------------
\section{Mechanizmy konsensusu w sieciach rozproszonych}

W środowisku rozproszonym, gdzie nie istnieje centralny arbiter decydujący o prawdziwości
danych, kluczowym wyzwaniem jest osiągnięcie porozumienia co do jednej, spójnej wersji
historii transakcji. Problem ten, znany w informatyce jako \textbf{problem bizantyjskich
generałów}, w sieciach blockchain rozwiązywany jest za pomocą mechanizmów konsensusu.
Są to zestawy reguł ekonomicznych i kryptograficznych, które umożliwiają węzłom sieci
wspólne zatwierdzanie nowych bloków danych bez konieczności wzajemnego zaufania.

\subsection{Proof of Work (PoW) -- Dowód pracy}

Mechanizm Proof of Work, wprowadzony przez Bitcoina, był pierwszą praktyczną implementacją
konsensusu w otwartej sieci peer-to-peer. Jego działanie opiera się na rywalizacji
obliczeniowej: uczestnicy sieci (górnicy) rozwiązują kosztowne obliczeniowo zadanie
kryptograficzne polegające na znalezieniu takiej wartości liczbowej (\textit{nonce}),
która po dołączeniu do nagłówka bloku da w wyniku skrót mniejszy od zadanego progu
trudności. Proces ten wymaga wielokrotnego próbkowania i zużycia znacznych ilości energii
elektrycznej.

Główną zaletą PoW jest udowodnione empirycznie wysokie bezpieczeństwo. Koszt ewentualnego
ataku jest wprost proporcjonalny do ceny sprzętu i energii potrzebnej do prześcignięcia
uczciwych uczestników sieci. Wadą jest niska wydajność transakcyjna oraz negatywny wpływ
środowiskowy związany z konsumpcją energii.

\subsection{Proof of Stake (PoS) -- Dowód stawki}

Proof of Stake stanowi odpowiedź na problemy skalowalności i energochłonności PoW.
W modelu PoS prawo do tworzenia nowych bloków nie zależy od mocy obliczeniowej, lecz od
ilości monet zdeponowanych przez uczestnika w sieci (tzw.~\textit{stake}). Im więcej monet
użytkownik zablokuje jako zabezpieczenie, tym większe prawdopodobieństwo, że zostanie
wybrany do zatwierdzenia kolejnego bloku i otrzyma związaną z tym nagrodę.

Mechanizm ten wprowadza ekonomiczną gwarancję uczciwości. W przypadku próby oszustwa
zdeponowany stake może zostać częściowo lub całkowicie skonfiskowany przez protokół
(tzw.~\textit{slashing}). Atak na sieć staje się przez to nieopłacalny finansowo,
ponieważ agresor ryzykuje utratę własnego kapitału.

% -----------------------------------------------------------------------
\section{Kryptograficzne podstawy bezpieczeństwa}

Bezpieczeństwo systemów opartych na blockchain wynika nie tylko z mechanizmu konsensusu,
ale także z zastosowania zaawansowanych narzędzi kryptograficznych. Dwie kluczowe klasy
algorytmów zapewniają integralność danych i autoryzację działań: funkcje skrótu oraz
kryptografia asymetryczna.

\subsection{Kryptograficzne funkcje skrótu}

Funkcja skrótu \( H \) jest deterministycznym algorytmem matematycznym, który dla
dowolnych danych wejściowych \( m \) generuje unikalny, krótki ciąg wyjściowy
\( h = H(m) \) o stałej długości, zwany skrótem lub hashem. W kontekście blockchain
kluczowe są trzy właściwości kryptograficznych funkcji skrótu:

\begin{itemize}
    \item \textbf{Jednokierunkowość:} Na podstawie wartości skrótu \( h \) niemożliwe
          jest obliczeniowe odtworzenie oryginalnych danych wejściowych \( m \).
    \item \textbf{Odporność na kolizje:} Praktycznie niemożliwe jest znalezienie dwóch
          różnych wiadomości \( m_1 \) i \( m_2 \), dla których \( H(m_1) = H(m_2) \).
          Algorytm SHA-256 generuje skrót o długości 256 bitów, co daje przestrzeń
          \( 2^{256} \) możliwych wartości -- prawdopodobieństwo przypadkowej kolizji
          jest zaniedbywalnie małe.
    \item \textbf{Efekt lawinowy:} Nawet minimalna modyfikacja danych wejściowych
          skutkuje drastyczną, nieprzewidywalną zmianą całego skrótu wyjściowego.
          Właściwość ta umożliwia natychmiastowe wykrycie próby fałszerstwa.
\end{itemize}

W architekturze opisywanego systemu skrót SHA-256 pełni rolę cyfrowego odcisku palca
dokumentu. To właśnie ten odcisk, a nie sam plik, zostaje zapisany w transakcji
blockchain.

\subsection{Kryptografia asymetryczna i podpisy cyfrowe (ECDSA)}

Tożsamość uczestników i autoryzacja działań w sieciach blockchain opierają się na
kryptografii klucza publicznego (asymetrycznej). Wykorzystywany jest standard ECDSA
(ang.~\textit{Elliptic Curve Digital Signature Algorithm}) oparty na krzywej eliptycznej
\textit{secp256k1}.

Każdy uczestnik sieci dysponuje parą matematycznie powiązanych kluczy:
\begin{itemize}
    \item \textbf{Klucz prywatny} (\(k\)): Losowo wygenerowana liczba, która musi być
          bezwzględnie chroniona. Stanowi dowód tożsamości i własności środków.
    \item \textbf{Klucz publiczny} (\(K\)): Generowany na podstawie klucza prywatnego
          przez operację na krzywej eliptycznej (\(K = k \cdot G\)). Może być swobodnie
          udostępniany.
    \item \textbf{Adres portfela:} Powstaje przez dwukrotne haszowanie klucza publicznego
          (SHA-256, następnie RIPEMD-160), co skraca jego długość i dodaje warstwę
          bezpieczeństwa.
\end{itemize}

Podpis cyfrowy tworzony jest kluczem prywatnym dla konkretnej wiadomości (np.~transakcji).
Każdy węzeł może zweryfikować autentyczność podpisu kluczem publicznym, nie poznając
przy tym klucza prywatnego. Mechanizm ten pozwala na jednoznaczne przypisanie
odpowiedzialności za daną transakcję do konkretnego adresu.

% -----------------------------------------------------------------------
\section{Architektura sieci Peer-to-Peer}

Publiczne sieci blockchain funkcjonują w oparciu o architekturę równorzędnych połączeń
(ang.~\textit{peer-to-peer}, P2P). Nie istnieje tu hierarchia klient-serwer w tradycyjnym
rozumieniu -- każdy węzeł łączy się z kilkoma innymi węzłami, tworząc gęstą, odporną na
uszkodzenia sieć połączeń. Gdy użytkownik inicjuje transakcję, informacja o niej jest
propagowana do wszystkich osiągalnych węzłów (ang.~\textit{gossip protocol}), aż dotrze
do walidatora, który włączy ją do nowego bloku.

Taka architektura gwarantuje wysoką odporność na cenzurę. Wyłączenie grupy węzłów w
konkretnej lokalizacji geograficznej nie przerywa działania całej sieci, o ile istnieje
choć jedna ścieżka połączeniowa do jej pozostałych uczestników.

% -----------------------------------------------------------------------
\section{Smart kontrakty}

Smart kontrakt (ang.~\textit{smart contract}) to program komputerowy przechowywany i
wykonywany bezpośrednio na blockchainie. Pojęcie to zostało wprowadzone przez Nicka Szabo
w latach 90. XX~wieku, jednak praktyczne zastosowanie na szeroką skalę umożliwiła dopiero
sieć Ethereum, zaprojektowana przez Vitalika Buterina~\cite{buterin2014ethereum}.

Smart kontrakt działa autonomicznie -- po wdrożeniu na blockchain jego kod jest niezmienny
i publicznie weryfikowalny, a wykonanie następuje automatycznie po spełnieniu zdefiniowanych
warunków, bez udziału żadnej strony trzeciej. Właściwości te czynią smart kontrakty
naturalnym narzędziem do budowania systemów wymagających zaufania między stronami, które
się nie znają.

W kontekście niniejszego projektu smart kontrakt pełni rolę niezmiennego rejestru
dokumentów: zapisuje kryptograficzne skróty plików wraz z metadanymi i udostępnia funkcję
weryfikacji, której wynik jest publicznie sprawdzalny przez każdego uczestnika sieci.

% -----------------------------------------------------------------------
\section{IPFS -- rozproszony system plików}

IPFS (ang.~\textit{InterPlanetary File System}) to protokół i sieć peer-to-peer służąca
do przechowywania i udostępniania danych w sposób rozproszony. W odróżnieniu od
tradycyjnego internetu, gdzie pliki są adresowane przez lokalizację serwera (URL), w IPFS
każdy plik jest identyfikowany przez swój hash -- tzw.~CID
(ang.~\textit{Content Identifier}).

Kluczową konsekwencją tego podejścia jest to, że \textbf{zawartość pliku jest nierozłącznie
związana z jego adresem}. Jeśli plik zostanie choćby minimalnie zmodyfikowany, jego CID
zmieni się całkowicie -- nie można więc po cichu podmienić pliku pod tym samym adresem.
Własność ta doskonale uzupełnia model weryfikacji oparty na hashach zapisanych
w blockchainie.




% ==================== ROZDZIAŁ 2: UŻYTE TECHNOLOGIE I METODY 
(ROZSZERZONY) ====================
\chapter{Użyte technologie i metody implementacji}


\section{Warstwa aplikacyjna -- Backend ASP.NET Core 8}

Logika biznesowa systemu obiegu dokumentów, obejmująca zarządzanie użytkownikami, rolami, uprawnieniami oraz procesami akceptacji dokumentów, zaimplementowana została w języku C\# z wykorzystaniem frameworka ASP.NET Core 8. Wybór tego stosu technologicznego podyktowany został jego wysoką wydajnością, dojrzałością ekosystemu oraz silnym typowaniem, które redukuje ryzyko błędów w fazie produkcyjnej.

\textbf{ASP.NET Core} zapewnia warstwę prezentacji w postaci zestawu kontrolerów API RESTful. Oznacza to, że backend nie generuje interfejsu graficznego (HTML), a jedynie udostępnia zestaw adresów URL (endpointów), za pomocą których zewnętrzna aplikacja kliencka (np. przeglądarkowa napisana w React lub Angular) może komunikować się z serwerem. Każdy endpoint odpowiada konkretnej akcji, np. `POST /api/documents` tworzy nowy dokument, a `GET /api/documents/{id}` pobiera jego metadane.


\textbf{Entity Framework Core (EF Core)} pełni funkcję mapowania obiektowo-relacyjnego (ORM). Umożliwia ono programiście operowanie na bazie danych za pomocą obiektów i kolekcji języka C\#, bez konieczności ręcznego pisania zapytań w języku SQL. EF Core tłumaczy operacje wykonane na obiektach (np. `context.Users.Add(nowyUzytkownik)`) na odpowiednie komendy SQL. Mechanizm migracji (`Add-Migration`, `Update-Database`) pozwala na utrzymywanie spójności schematu relacyjnej bazy danych PostgreSQL z modelem danych aplikacji, umożliwiając łatwe wprowadzanie zmian w strukturze projektu w trakcie jego rozwoju.

\textbf{PostgreSQL} został wybrany jako główny system zarządzania bazą danych (RDBMS). W przeciwieństwie do lekkich baz plikowych (np. SQLite), PostgreSQL oferuje pełne wsparcie dla transakcji ACID, zaawansowane mechanizmy integralności referencyjnej oraz wysoką skalowalność. W kontekście systemu obiegu dokumentów przechowuje on dane, które nie wymagają niezmienności oferowanej przez blockchain, ale muszą być szybko wyszukiwane i modyfikowane -- są to przede wszystkim profile użytkowników, uprawnienia, historia akceptacji oraz metadane plików (nazwy, opisy, odnośniki do IPFS).

\newpage
\section{Bezpieczeństwo komunikacji i przechowywania danych uwierzytelniających}

W systemie, który z założenia operuje na wrażliwych dokumentach, kluczowym aspektem jest bezpieczeństwo warstwy dostępu. Zastosowano dwa główne mechanizmy ochrony.

\subsection{Uwierzytelnianie za pomocą tokenów JWT}
Mechanizm sesji serwerowych opartych na ciasteczkach został zastąpiony przez nowoczesny standard JSON Web Token (JWT). Po pomyślnym zalogowaniu (poprzez podanie loginu i hasła lub podpisu kryptograficznego portfela DIVI), backend generuje podpisany cyfrowo token zawierający zakodowane informacje o użytkowniku (np. identyfikator, rolę) oraz czas ważności. Token ten jest odsyłany do klienta, który musi dołączać go do nagłówka `Authorization` każdego kolejnego żądania HTTP.
Zastosowanie JWT zapewnia bezstanowość (ang. \textit{statelessness}) serwera. Oznacza to, że serwer nie musi przechowywać informacji o zalogowanych użytkownikach w pamięci RAM ani w bazie danych -- weryfikacja tożsamości odbywa się wyłącznie na podstawie kryptograficznej poprawności podpisu tokena. Znacząco ułatwia to skalowanie aplikacji w środowisku chmurowym oraz redukuje powierzchnię ataku związaną z kradzieżą sesji.

\subsection{Bezpieczne przechowywanie haseł -- Argon2id}
Bezpieczeństwo tradycyjnego loginu i hasła opiera się na algorytmie mieszającym Argon2id, który został zwycięzcą konkursu Password Hashing Competition (PHC) i jest rekomendowany przez OWASP jako najlepsza praktyka. W przeciwieństwie do starszych algorytmów (jak MD5 czy SHA-1), Argon2id został zaprojektowany tak, aby był odporny na ataki z wykorzystaniem dedykowanych układów scalonych (ASIC) oraz kart graficznych (GPU). Dzięki konfigurowalnym parametrom kosztu pamięci, czasu i równoległości, proces weryfikacji hasła jest celowo spowolniony i wymaga znacznych zasobów RAM, co czyni ataki słownikowe typu \textit{brute-force} ekonomicznie nieopłacalnymi dla potencjalnego włamywacza.

\newpage
\section{Warstwa integracji z blockchainem -- Ethereum, Solidity i Hardhat}

Chociaż logika aplikacji osadzona jest w ekosystemie DIVI, rejestracja dowodów istnienia dokumentów realizowana jest na kompatybilnej maszynie wirtualnej Ethereum (EVM). Takie podejście umożliwia wykorzystanie dojrzałych narzędzi deweloperskich Ethereum i potencjalną migrację do innych sieci EVM w przyszłości.

\subsection{Smart Kontrakt w języku Solidity}
Smart kontrakt stanowi samowykonalny program komputerowy, którego kod i stan są przechowywane bezpośrednio w łańcuchu bloków. W projekcie napisano kontrakt w języku Solidity, który pełni rolę zdecentralizowanego rejestru metadanych dokumentów. Jego główne funkcje publiczne to:
\begin{itemize}
    \item \texttt{registerDocument(bytes32 fileHash, string memory ipfsCid, string memory metadata)} -- pozwala autoryzowanemu użytkownikowi (lub kontraktowi backendowemu) na zapisanie trwałego wpisu w blockchainie. Funkcja ta emituje zdarzenie (ang. \textit{event}) `DocumentRegistered`, które może być nasłuchiwane przez aplikację frontendową w czasie rzeczywistym.
    \item \texttt{verifyDocument(bytes32 fileHash)} -- metoda widoku (nie modyfikująca stanu), która pozwala każdemu (nawet osobom nieposiadającym konta w systemie) zweryfikować, czy dany skrót pliku został wcześniej zarejestrowany i przez kogo.
\end{itemize}
Należy podkreślić, że kontrakt nie przechowuje całego dokumentu, a jedynie jego unikalny odcisk (hash SHA-256) oraz odnośnik do lokalizacji pliku w sieci IPFS. Dzięki temu koszt transakcji (gas) utrzymywany jest na akceptowalnym poziomie.

\subsection{Środowisko deweloperskie Hardhat}
Hardhat to kompleksowe środowisko programistyczne dla projektów Ethereum. W kontekście niniejszej pracy zapewnia ono powtarzalność procesu wdrażania i testowania.
\begin{itemize}
    \item \textbf{Kompilacja:} Narzędzie kompiluje kod źródłowy Solidity do kodu bajtowego EVM.
    \item \textbf{Lokalny blockchain:} Uruchamia na komputerze dewelopera tymczasową, prywatną sieć Ethereum, na której można testować działanie kontraktu bez ponoszenia kosztów transakcji.
    \item \textbf{Testy automatyczne:} Skrypty testowe napisane w JavaScript (z wykorzystaniem bibliotek Mocha i Chai) symulują interakcje z kontraktem, sprawdzając poprawność rejestracji, mechanizmy kontroli dostępu oraz odporność na próby fałszerstwa.
    \item \textbf{Wdrażanie (Deployment):} Hardhat automatyzuje proces wysyłania transakcji tworzącej kontrakt na sieć testową Sepolia, zapisując adres wdrożonego kontraktu do późniejszego wykorzystania przez backend C\#.
\end{itemize}

\section{Rozproszone przechowywanie plików -- IPFS i Pinata}

Bezpośrednie przechowywanie plików w bazie danych PostgreSQL lub w transakcjach blockchain jest nieefektywne kosztowo i technicznie (ograniczenia rozmiaru bloku, przepustowość sieci). W związku z tym system wykorzystuje protokół IPFS (InterPlanetary File System).

\textbf{IPFS} to sieć P2P zaprojektowana do przechowywania i udostępniania danych w sposób rozproszony. Kluczową cechą IPFS jest adresowanie zawartości (ang. \textit{content addressing}). W tradycyjnym Internecie (HTTP) plik znajduje się pod konkretnym adresem URL (np. \texttt{serwer.pl/dokument.pdf}). Jeśli plik zostanie przeniesiony lub serwer przestanie działać, link przestaje działać. W IPFS identyfikatorem pliku jest jego kryptograficzny skrót (CID). Oznacza to, że:
\begin{enumerate}
    \item Plik jest weryfikowalny -- klient pobierający plik może sam obliczyć jego hash i porównać z CID, upewniając się, że dane nie zostały uszkodzone ani sfałszowane.
    \item Plik jest odporny na cenzurę -- tak długo, jak przynajmniej jeden węzeł w sieci IPFS przechowuje kopię pliku, jest on dostępny pod tym samym CID.
\end{enumerate}

W projekcie wykorzystano usługę \textbf{Pinata} jako bramę IPFS (ang. \textit{pinning service}). Usługa ta dostarcza wygodne REST API, za pomocą którego backend C\# wysyła plik dokumentu. Pinata przechowuje plik na swoich serwerach węzłów IPFS oraz "przypina" go (ang. \textit{pin}), co gwarantuje, że plik nie zostanie usunięty przez mechanizm odśmiecania sieci (ang. \textit{garbage collection}). W odpowiedzi zwracany jest identyfikator CID, który następnie zostaje zapisany w smart kontrakcie jako trwały dowód powiązania dokumentu z plikiem.

\section{Sieć testowa Sepolia jako środowisko uruchomieniowe}

Całość operacji blockchainowych projektu realizowana jest na publicznej sieci testowej Ethereum o nazwie Sepolia. Wybór ten motywowany jest kilkoma czynnikami:
\begin{itemize}
    \item \textbf{Brak kosztów rzeczywistych:} Waluta Ether na Sepolii (SepoliaETH) nie posiada wartości rynkowej i jest dystrybuowana bezpłatnie przez tzw. krany (ang. \textit{faucets}). Pozwala to na swobodne testowanie wszystkich ścieżek działania systemu, w tym błędnych, bez ponoszenia strat finansowych.
    \item \textbf{Identyczność techniczna:} Sieć Sepolia pod względem parametrów bloków, mechanizmu konsensusu (Proof of Stake) i działania maszyny wirtualnej EVM jest wierną kopią głównej sieci Ethereum (Mainnet). Gwarantuje to, że kod działający na testnecie będzie działał identycznie po ewentualnej migracji na produkcję.
\end{itemize}










% ==================== ROZDZIAŁ 3: REALIZACJA PRAKTYCZNA SYSTEMU (NOWY ROZDZIAŁ) ====================
\chapter{Realizacja praktyczna systemu obiegu dokumentów}

\section{Architektura logiczna systemu}

System został zaprojektowany jako aplikacja hybrydowa, łącząca scentralizowaną logikę
biznesową (backend C\#) z zdecentralizowaną warstwą zaufania (blockchain oraz IPFS).
Podejście to pozwala na zachowanie wygody tradycyjnych aplikacji webowych przy jednoczesnym
zapewnieniu kryptograficznie gwarantowanej niezmienności kluczowych danych. Można wyróżnić
cztery główne komponenty współpracujące ze sobą w sposób asynchroniczny:

\begin{enumerate}
    \item \textbf{Klient (Frontend Web App):} Interfejs użytkownika działający w przeglądarce
          internetowej. Odpowiada za interakcję z użytkownikiem: formularze logowania, listy
          dokumentów, upload plików, moduł weryfikatora oraz panel administracyjny.
          Komunikuje się wyłącznie z Backendem API poprzez żądania HTTP.

    \item \textbf{Backend API (ASP.NET Core 8):} Centralny punkt koordynujący całą logikę
          biznesową systemu. Odpowiada za autoryzację żądań, zarządzanie stanem procesu
          akceptacji dokumentów w bazie PostgreSQL oraz orkiestrację komunikacji
          z zewnętrznymi usługami (Pinata, sieć Ethereum Sepolia).

    \item \textbf{Warstwa Blockchain (Smart Kontrakt na Sepolii):} Pełni rolę
          niezmienialnego notariusza. Przechowuje dowód istnienia dokumentu (hash SHA-256
          oraz CID IPFS) wraz z podstawowymi metadanymi (typ dokumentu, adres wystawcy,
          data rejestracji on-chain). Żaden podmiot -- w tym twórcy systemu -- nie jest
          w stanie zmodyfikować ani usunąć raz zapisanych danych.

    \item \textbf{Warstwa Przechowywania Plików (IPFS + Pinata):} Odpowiada za fizyczne
          przechowywanie treści plików w sposób rozproszony i odporny na utratę danych.
          Dostęp do pliku możliwy jest przez dowolną publiczną bramę IPFS wyłącznie
          na podstawie jego CID, bez zależności od konkretnego serwera.
\end{enumerate}

Całość systemu zaprojektowano zgodnie z zasadą rozdzielenia odpowiedzialności
(ang.~\textit{Separation of Concerns}): baza PostgreSQL przechowuje dane operacyjne
i aplikacyjne, blockchain przechowuje niezmienne dowody integralności, natomiast IPFS
przechowuje rzeczywistą treść plików. Żadna z tych warstw nie jest zbędna -- ich
współpraca tworzy kompletny i odporny system.

% -----------------------------------------------------------------------
\newpage
\section{Przepływ procesu rejestracji dokumentu}

Opisany poniżej scenariusz ilustruje ścieżkę danych od momentu wyboru pliku przez
użytkownika do ostatecznego potwierdzenia jego autentyczności w łańcuchu bloków. Proces
zaprojektowano tak, aby był transparentny dla użytkownika końcowego, który obserwuje
jedynie standardowy wskaźnik postępu operacji.

\begin{enumerate}
    \item \textbf{Inicjalizacja:} Użytkownik, po pomyślnym uwierzytelnieniu tokenem JWT,
          wybiera plik z dysku lokalnego i wypełnia formularz metadanymi: nazwa dokumentu,
          opis, kategoria, lista odbiorców oraz data ważności (opcjonalna).

    \item \textbf{Przetwarzanie wstępne w Backendzie:}
    \begin{itemize}
        \item Backend oblicza skrót SHA-256 odebranego strumienia bajtów pliku. Obliczenie
              odbywa się w pamięci serwera bez zapisywania pliku na dysku lokalnym.
        \item Przeprowadzana jest walidacja -- sprawdzenie, czy dokument o tym samym hashu
              nie został już zarejestrowany w bazie PostgreSQL (zapobieganie duplikacji
              wpisów lokalnych) oraz czy format i rozmiar pliku spełniają wymagania systemu.
        \item Weryfikowany jest status wystawcy -- tylko użytkownicy z rolą zatwierdzonego
              wystawcy (ang.~\textit{issuer}) mogą inicjować rejestrację nowego dokumentu.
    \end{itemize}

    \item \textbf{Zapis w IPFS:} Backend nawiązuje połączenie z API usługi Pinata
          i przesyła plik. W odpowiedzi otrzymuje unikalny identyfikator treści (CID),
          który od tej chwili jednoznacznie identyfikuje ten plik w sieci IPFS niezależnie
          od jakiegokolwiek serwera.

    \item \textbf{Zapis w Blockchain:} Backend, wykorzystując bibliotekę Nethereum oraz
          klucz prywatny dedykowanego portfela systemowego (tzw.~\textit{relay wallet}),
          konstruuje transakcję wywołującą metodę \texttt{registerDocument} w smart
          kontrakcie. Transakcja zawiera hash SHA-256 oraz CID dokumentu. Po
          rozpropagowaniu transakcji w sieci Sepolia i włączeniu jej do bloku dane stają
          się publicznie niezmienialnie zapisane w historii łańcucha.

    \item \textbf{Zapis lokalny i powiadomienie:} Po otrzymaniu potwierdzenia z sieci
          (numer bloku i hash transakcji), Backend zapisuje w bazie PostgreSQL pełny rekord
          dokumentu -- włączając CID, hash SHA-256, identyfikator wystawcy, metadane oraz
          początkowy status dokumentu -- i zwraca odpowiedź sukcesu do klienta. Wystawca
          otrzymuje potwierdzenie z linkiem do transakcji w eksploratorze bloków Sepolia.
\end{enumerate}

% -----------------------------------------------------------------------
\newpage
\section{Proces weryfikacji autentyczności dokumentu}

Jedną z kluczowych funkcjonalności systemu jest możliwość niezależnej weryfikacji, czy
dany plik odpowiada zapisowi znajdującemu się w blockchainie. Proces ten nie wymaga
logowania ani posiadania specjalnych uprawnień -- jest dostępny publicznie.

\begin{enumerate}
    \item Użytkownik zewnętrzny (np.~audytor, odbiorca dokumentu) otrzymuje plik dokumentu
          dowolnym kanałem komunikacji.
    \item Korzystając z publicznego modułu weryfikatora w aplikacji, wgrywa plik do
          dedykowanego formularza.
    \item Aplikacja frontendowa oblicza lokalnie skrót SHA-256 pliku bezpośrednio
          w przeglądarce użytkownika, bez przesyłania pliku na serwer. Rozwiązanie to
          chroni prywatność treści dokumentu -- serwer nigdy nie widzi jego zawartości.
    \item Obliczony skrót wysyłany jest do Backendu, który wywołuje funkcję
          \texttt{verifyDocument} w smart kontrakcie na sieci Sepolia.
    \item Smart kontrakt przeszukuje swoją wewnętrzną mapę rejestrów i zwraca wynik.
    \item W przypadku znalezienia pasującego hasha zwracane są dane rejestracji: adres
          portfela wystawcy, numer bloku i data rejestracji, CID pliku w IPFS oraz
          typ dokumentu.
    \item Backend, na podstawie CID, pobiera plik referencyjny z IPFS i porównuje jego
          zawartość z plikiem dostarczonym do weryfikacji. Zgodność stanowi ostateczny,
          matematyczny dowód autentyczności dokumentu.
    \item W przypadku braku zgodności lub nieznalezienia hasha w kontrakcie system zwraca
          jednoznaczną informację o negatywnym wyniku weryfikacji, wskazując czy dokument
          jest nieznany systemowi, czy też mógł zostać zmodyfikowany po rejestracji.
\end{enumerate}

% -----------------------------------------------------------------------
\newpage
\section{Przeglądanie i zarządzanie dokumentami}

\subsection{Lista dokumentów wystawcy}

Każdy zalogowany wystawca ma dostęp do panelu z pełną historią wystawionych przez siebie
dokumentów. Lista prezentuje dla każdej pozycji: nazwę dokumentu, datę rejestracji,
aktualny status (oczekujący, zaakceptowany, odrzucony, wygasły), hash SHA-256 skrócony
do formy czytelnego identyfikatora oraz link do transakcji rejestrującej w eksploratorze
bloków Sepolia. Lista obsługuje stronicowanie oraz filtrowanie według statusu i zakresu
dat, co umożliwia sprawne zarządzanie dużą liczbą dokumentów.

\subsection{Podgląd szczegółów dokumentu}

Po wybraniu konkretnego dokumentu z listy wyświetlany jest widok szczegółowy zawierający
pełne metadane: kompletny hash SHA-256 w formacie szesnastkowym, adres CID z aktywnym
linkiem do publicznej bramy IPFS, adres portfela Ethereum, który dokonał rejestracji,
numer bloku, hash transakcji blockchain oraz listę odbiorców powiązanych z dokumentem.
Z poziomu tego widoku możliwe jest pobranie oryginalnego pliku z IPFS oraz wygenerowanie
raportu weryfikacyjnego w formacie PDF zawierającego wszystkie dane potwierdzające
autentyczność dokumentu.

\subsection{Historia operacji na dokumencie}

Dla każdego dokumentu prowadzony jest dziennik zdarzeń (ang.~\textit{audit log})
rejestrowany w bazie PostgreSQL. Dziennik obejmuje: moment przesłania dokumentu do
systemu, zmianę statusu przez administratora, każdorazowe pobranie pliku przez odbiorcę
oraz próby weryfikacji (zarówno zakończone powodzeniem, jak i nieudane). Dziennik ten
nie jest modyfikowalny przez zwykłych użytkowników i stanowi uzupełnienie dowodów
zapisanych na blockchainie w zakresie danych operacyjnych.

% -----------------------------------------------------------------------
\section{Mechanizm wygasania i odwoływania dokumentów}

System obsługuje opcjonalne daty ważności dokumentów. Wystawca może podczas rejestracji
określić datę, po której dokument powinien zostać uznany za nieaktywny. Po upływie tego
terminu status dokumentu w bazie PostgreSQL jest automatycznie aktualizowany przez
zadanie cykliczne działające w tle (ang.~\textit{background job}).

Ważne jest podkreślenie, że wygaśnięcie dokumentu w warstwie aplikacyjnej nie usuwa
danych z blockchainu -- hash i CID pozostają tam na zawsze. Wygaśnięcie oznacza jedynie,
że system nie będzie traktował dokumentu jako aktualnie obowiązującego przy weryfikacji
statusu biznesowego, co jest adekwatne np.~dla dokumentów takich jak umowy, certyfikaty
czy pełnomocnictwa o ograniczonym czasie obowiązywania. Proces weryfikacji matematycznej
(porównanie hasha z łańcuchem) nadal zwraca wynik pozytywny, co potwierdza, że plik nie
był modyfikowany -- informacja o wygaśnięciu jest zwracana jako osobne pole w odpowiedzi.

% -----------------------------------------------------------------------
\section{Zarządzanie tożsamością i rolami w systemie}

System wprowadza podział na role użytkowników w celu odwzorowania rzeczywistych procesów
biznesowych w organizacji. Mechanizm ten realizowany jest w warstwie backendowej C\#
i bazie danych PostgreSQL.

Wyróżnione zostały następujące role:
\begin{itemize}
    \item \textbf{Gość (Guest):} Użytkownik niezalogowany, który ma dostęp wyłącznie
          do publicznego modułu weryfikatora dokumentów.
    \item \textbf{Użytkownik (User):} Konto zarejestrowane w systemie. Może przeglądać
          dokumenty, do których zostało mu przyznane prawo dostępu jako odbiorcy,
          oraz składać wnioski o nadanie roli wystawcy.
    \item \textbf{Wystawca (Issuer):} Rola przyznawana przez administratora po pozytywnym
          rozpatrzeniu wniosku. Umożliwia rejestrację nowych dokumentów w blockchain
          oraz zarządzanie własnymi dokumentami.
    \item \textbf{Administrator (Admin):} Pełny dostęp do systemu. Odpowiada za
          rozpatrywanie wniosków o rolę wystawcy, zarządzanie kontami użytkowników,
          podgląd dzienników systemowych oraz konfigurację parametrów aplikacji.
\end{itemize}

Każde żądanie do API jest weryfikowane pod kątem posiadanego tokenu JWT oraz przypisanej
do konta roli. Próba wykonania operacji bez odpowiednich uprawnień skutkuje odpowiedzią
HTTP 403 (Forbidden) i wpisem do dziennika zdarzeń bezpieczeństwa.

% -----------------------------------------------------------------------
\section{Proces wnioskowania o rolę wystawcy}

Nadawanie uprawnień wystawcy nie odbywa się automatycznie -- wymaga złożenia formalnego
wniosku i jego zatwierdzenia przez administratora. Proces ten odwzorowuje realne wymogi
organizacyjne, gdzie nie każdy użytkownik systemu powinien mieć prawo do rejestrowania
dokumentów w imieniu organizacji.

Użytkownik wypełnia formularz wniosku, podając uzasadnienie oraz -- opcjonalnie --
adres portfela Ethereum, którego zamierza używać do podpisywania transakcji. Wniosek
trafia do kolejki w panelu administracyjnym, gdzie administrator może go zatwierdzić
lub odrzucić z podaniem powodu. Po zatwierdzeniu rola użytkownika jest aktualizowana
w bazie, a on sam otrzymuje powiadomienie. Cała historia wniosków jest archiwizowana
w celach audytowych.

% -----------------------------------------------------------------------
\newpage
\section{Uwierzytelnianie z użyciem portfela Ethereum}

Oprócz standardowego logowania przez e-mail i hasło, system udostępnia mechanizm
uwierzytelniania przy użyciu portfela Ethereum (np.~MetaMask). Proces ten opiera się
na kryptograficznym dowodzie posiadania klucza prywatnego, bez przesyłania żadnych
poświadczeń:

\begin{enumerate}
    \item Użytkownik inicjuje logowanie portfelem, podając adres swojego konta Ethereum.
    \item Backend generuje jednorazowy losowy ciąg znaków (nonce) powiązany z tym adresem
          i odsyła go do klienta.
    \item Użytkownik podpisuje nonce swoim portfelem (operacja wykonywana lokalnie,
          klucz prywatny nigdy nie opuszcza urządzenia).
    \item Podpis oraz adres portfela są wysyłane do Backendu.
    \item Backend weryfikuje, czy podpis został rzeczywiście utworzony kluczem prywatnym
          odpowiadającym podanemu adresowi. Jeśli weryfikacja przebiegnie pomyślnie,
          użytkownik otrzymuje token JWT ważny przez określony czas.
    \item Nonce jest natychmiast unieważniany po użyciu, co uniemożliwia jego ponowne
          wykorzystanie w ewentualnym ataku odtworzeniowym (ang.~\textit{replay attack}).
\end{enumerate}

% -----------------------------------------------------------------------
\section{Bezpieczeństwo systemu}

Bezpieczeństwo systemu zapewniane jest na kilku niezależnych poziomach, co zgodne jest
z zasadą obrony warstwowej (ang.~\textit{defense in depth}).

Hasła użytkowników nigdy nie są przechowywane w postaci jawnej. Do ich haszowania
stosowany jest algorytm \textbf{Argon2id}, uznawany za aktualny standard w tej dziedzinie.
Algorytm ten jest celowo kosztowny obliczeniowo i pamięciowo, co drastycznie spowalnia
potencjalne ataki słownikowe i brute-force.

Komunikacja między wszystkimi komponentami systemu odbywa się wyłącznie przez szyfrowane
połączenia HTTPS/TLS. Wrażliwe dane konfiguracyjne (klucze prywatne portfela systemowego,
klucze API Pinata) przechowywane są jako zmienne środowiskowe, nigdy w kodzie źródłowym
ani repozytorium.

\newpage
Dostęp do poszczególnych endpointów API jest ograniczany atrybutami ról bezpośrednio
w kodzie kontrolerów C\#, co eliminuje ryzyko przypadkowego pominięcia weryfikacji
uprawnień. Każda próba nieautoryzowanego dostępu jest rejestrowana w dzienniku zdarzeń
bezpieczeństwa z adresem IP i identyfikatorem żądania, co umożliwia późniejszą analizę
i detekcję prób ataku.

\begin{itemize}
    \item \textbf{Użytkownik (User):} Podstawowa rola umożliwiająca tworzenie nowych dokumentów, przeglądanie dokumentów własnych oraz tych udostępnionych.
    \item \textbf{Menedżer (Manager):} Rola posiadająca uprawnienia do zatwierdzania (akceptacji) lub odrzucania dokumentów w ramach zdefiniowanego obiegu. Menedżer może również przeglądać dokumenty podległych mu działów.
    \item \textbf{Administrator (Admin):} Rola nadrzędna z dostępem do panelu zarządzania użytkownikami, definiowania struktur organizacyjnych oraz konfiguracji parametrów systemu (np. adresu smart kontraktu).
\end{itemize}
Mapowanie pomiędzy użytkownikiem systemu a jego adresem w sieci DIVI (lub Ethereum) jest opcjonalne i wykorzystywane głównie do potwierdzenia tożsamości poprzez podpisanie wiadomości kluczem prywatnym, co stanowi dodatkowy, silny czynnik uwierzytelniania.


\begin{thebibliography}{9}
\bibitem{nakamoto2008bitcoin}
S.~Nakamoto, \textit{Bitcoin: A Peer-to-Peer Electronic Cash System}, 2008. [Online]. Dostęp: \url{https://bitcoin.org/bitcoin.pdf}

\bibitem{ipfs}
J.~Benet, \textit{IPFS - Content Addressed, Versioned, P2P File System}, 2014. [Online]. Dostęp: \url{https://ipfs.io/ipfs/QmR7GSQM93Cx5eAg6a6yRzNde1FQv7uL6X1o4k7zrJa3LX/ipfs.draft3.pdf}

\bibitem{ethereum}
V.~Buterin, \textit{Ethereum: A Next-Generation Smart Contract and Decentralized Application Platform}, 2014. [Online]. Dostęp: \url{https://ethereum.org/en/whitepaper/}
\end{thebibliography}

\end{document}

\end{document}